% ============================================================
% DOCUMENTCLASS
% ============================================================
\documentclass[12pt,oneside]{book}

% ============================================================
% METADATOS
% ============================================================
\newcommand{\universidad}{Universidad de Buenos Aires (UBA)}
\newcommand{\facultad}{Facultad de Derecho}
\newcommand{\materia}{Derecho Procesal Civil}
\newcommand{\titulo}{Práctico Profesional: notificaciones, demanda, prueba y alegatos en el proceso civil}
\newcommand{\autor}{Isaac Edgar Camacho Ocampo}
\newcommand{\anio}{2026}

% ============================================================
% PAQUETES
% ============================================================
\usepackage[utf8]{inputenc}
\usepackage[T1]{fontenc}
\usepackage[spanish,es-nodecimaldot]{babel}

\usepackage[
    top=2.5cm,
    bottom=2.5cm,
    left=3cm,
    right=2.5cm,
    headheight=15pt
]{geometry}

\usepackage{amsmath}
\usepackage{amssymb}
\usepackage{mathtools}

\usepackage{graphicx}
\usepackage{float}
\usepackage{caption}
\usepackage{subcaption}

\usepackage{booktabs}
\usepackage{array}
\usepackage{longtable}
\usepackage{tabularx}

\usepackage{enumitem}

\usepackage{xcolor}
\usepackage[most]{tcolorbox}

\usepackage{fancyhdr}
\usepackage{ragged2e}

\usepackage[
    colorlinks=true,
    linkcolor=blue!50!black,
    citecolor=green!40!black,
    urlcolor=blue!60!black,
    pdftitle={\titulo},
    pdfauthor={\autor},
    pdfsubject={Apuntes universitarios},
    bookmarks=true
]{hyperref}

% ============================================================
% CONFIGURACIÓN
% ============================================================
\graphicspath{
    {./img/}
    {./graficos/}
    {./figuras/}
}

\setcounter{tocdepth}{2}

\setlength{\parindent}{0pt}
\setlength{\parskip}{0.5em}

\setlist[itemize]{
    topsep=0.3em,
    itemsep=0.2em
}

\setlist[enumerate]{
    topsep=0.3em,
    itemsep=0.2em
}

\clubpenalty=10000
\widowpenalty=10000

% ============================================================
% COMANDOS
% ============================================================
\newcommand{\abs}[1]{\left|#1\right|}
\newcommand{\norm}[1]{\left\lVert#1\right\rVert}
\newcommand{\vect}[1]{\mathbf{#1}}
\newcommand{\deriv}[2]{\frac{d#1}{d#2}}
\newcommand{\pderiv}[2]{\frac{\partial#1}{\partial#2}}

\newcolumntype{Y}{>{\raggedright\arraybackslash}X}
\newcolumntype{L}[1]{>{\raggedright\arraybackslash}p{#1}}

% ============================================================
% ENTORNOS / CAJAS
% ============================================================
\newtcolorbox{definition}[1]{
    enhanced,
    breakable,
    title={Definición: #1},
    fonttitle=\bfseries,
    colback=blue!3,
    colframe=blue!50!black,
    boxrule=0.8pt,
    arc=2mm
}

\newtcolorbox{important}{
    enhanced,
    breakable,
    title={Importante},
    fonttitle=\bfseries,
    colback=yellow!8,
    colframe=orange!70!black,
    boxrule=0.8pt,
    arc=2mm
}

\newtcolorbox{example}[1]{
    enhanced,
    breakable,
    title={Ejemplo: #1},
    fonttitle=\bfseries,
    colback=green!4,
    colframe=green!40!black,
    boxrule=0.8pt,
    arc=2mm
}

\newtcolorbox{formula}[1]{
    enhanced,
    breakable,
    title={Fórmula / Regla: #1},
    fonttitle=\bfseries,
    colback=purple!4,
    colframe=purple!50!black,
    boxrule=0.8pt,
    arc=2mm
}

\newtcolorbox{exercise}[1]{
    enhanced,
    breakable,
    title={Ejercicio: #1},
    fonttitle=\bfseries,
    colback=gray!5,
    colframe=gray!60!black,
    boxrule=0.8pt,
    arc=2mm
}

\newtcolorbox{note}{
    enhanced,
    breakable,
    title={Nota},
    fonttitle=\bfseries,
    colback=white,
    colframe=gray!50,
    boxrule=0.6pt,
    arc=2mm
}

% ============================================================
% CONFIGURACIÓN FINAL
% ============================================================
\pagestyle{fancy}
\fancyhf{}
\fancyhead[L]{\nouppercase{\leftmark}}
\fancyhead[R]{\materia}
\fancyfoot[C]{\thepage}
\renewcommand{\headrulewidth}{0.4pt}
\renewcommand{\footrulewidth}{0pt}

\fancypagestyle{plain}{
    \fancyhf{}
    \fancyfoot[C]{\thepage}
    \renewcommand{\headrulewidth}{0pt}
}

% ============================================================
% DOCUMENT
% ============================================================
\begin{document}

% ------------------------------------------------------------
% PORTADA
% ------------------------------------------------------------
\begin{titlepage}

\thispagestyle{empty}

\begin{center}

\vspace*{1cm}

{\large \textsc{\universidad}}

\vspace{0.2cm}

{\large \textsc{\facultad}}

\vspace{3cm}

{\Huge \textbf{\materia}}

\vspace{1cm}

\parbox{0.85\textwidth}{\centering {\LARGE \titulo}}

\vspace{0.8cm}

\textit{Comprender los conceptos es el primer paso para convertir el estudio en conocimiento.}

\vspace{1.2cm}

\rule{0.70\textwidth}{0.5pt}

\vspace{0.5cm}

{\large Apuntes de estudio}

\vspace{0.5cm}

\rule{0.70\textwidth}{0.5pt}

\vfill

{\large \autor}

\vspace{0.3cm}

{\large \anio}

\end{center}

\vfill

\begin{flushleft}
\textit{Este es un modesto aporte a los estudiantes de ``\materia'' no pretende sustituir el material oficial de la materia.}
\end{flushleft}

\end{titlepage}

% ------------------------------------------------------------
% ÍNDICE
% ------------------------------------------------------------
\tableofcontents

% ------------------------------------------------------------
% CONTENIDO
% ------------------------------------------------------------

% ============================================================
\chapter[Presentación de la clase]{Presentación de la clase}
% ============================================================

\section{¿Qué se trató en esta clase?}

La clase corresponde al 10 de noviembre de 2026 y sigue el recorrido de un juicio civil desde que se presenta la demanda hasta los alegatos, con una pregunta práctica detrás de cada tema: \textbf{¿cómo se hace, en qué plazo y qué puede salir mal?}

Los contenidos se agrupan en cuatro grandes bloques:

\begin{itemize}
    \item \textbf{Notificaciones (arts.\ 135 y 136 del Código Procesal).} Qué resoluciones van siempre por cédula (la demanda y la sentencia), qué medios admite la ley, cuándo corresponde un edicto, cómo se tramitan las cédulas (ley, papel y electrónica) y cuándo se puede usar carta documento, WhatsApp o mail.
    \item \textbf{Notificación por nota, acordadas y plazos.} Los días de nota (martes y viernes), qué es una acordada, cómo se «deja nota» en el PJN y por qué, ante la duda, hay que contar corto.
    \item \textbf{Demanda y primer proveído.} Cómo se arma la demanda (art.\ 330), cómo se interpreta «notifíquese», «hágase saber» o «traslado», y hasta cuándo se puede ampliar la demanda o la prueba.
    \item \textbf{Prueba y alegatos.} Cómo se ofrece y fundamenta la prueba, la audiencia 360 y los 40 días hábiles, la diferencia entre negligencia o caducidad y prueba denegada, y cómo se estructura el alegato.
\end{itemize}

\section{¿Para qué sirve esto en la vida profesional?}

\begin{important}
\textbf{Sirve para no perder un caso por un error de forma.} Un abogado puede tener razón en el fondo y perder por una notificación nula, un plazo mal contado o una prueba mal ofrecida. Lo que se vio en esta clase es el \textbf{trabajo diario del litigante}: confeccionar cédulas, controlar proveídos, dejar nota, contar plazos, pedir pruebas y alegar.
\end{important}

\begin{itemize}
    \item \textbf{Evita nulidades y caducidades.} Una cédula mal dirigida o un plazo mal contado puede hacer perder derechos.
    \item \textbf{Ahorra tiempo y dinero.} Saber ir directo a la oficina de notificaciones o consultar el código de barras evita demoras de semanas.
    \item \textbf{Se usa desde el primer día.} Los patrocinios gratuitos del práctico y cualquier estudio jurídico exigen manejar estas herramientas.
\end{itemize}

\section{Utilidad para la vida personal}

Los contenidos también resultan útiles fuera del ejercicio profesional, porque cualquier persona puede recibir una cédula, una carta documento o un reclamo, y lo aprendido enseña a reaccionar correctamente:

\begin{itemize}
    \item \textbf{Leer una notificación y actuar a tiempo:} saber que existen plazos y que pueden ser de pocas horas o pocos días.
    \item \textbf{Documentar:} guardar constancias (códigos, comprobantes, capturas) que permiten probar qué pasó y cuándo.
    \item \textbf{Reclamar con método:} explicar con claridad qué se pide y por qué, ofrecer pruebas y resumir los argumentos, tanto en un juicio como en un reclamo a una empresa.
    \item \textbf{Pensar en riesgos:} aplicar el criterio «ante la duda, contar corto» a cualquier compromiso con fecha límite.
\end{itemize}

\section{Relación entre los temas: un ejemplo cotidiano}

Los bloques de la clase pueden compararse con los pasos de un reclamo por la compra online de un producto defectuoso:

\begin{example}{reclamo por un producto defectuoso comprado online}
\begin{itemize}
    \item \textbf{Bloque 1 (qué se notifica y por qué medio):} la tienda tiene que enterarse de forma oficial; no alcanza con avisarle por cualquier lado.
    \item \textbf{Bloque 2 (la cédula en la práctica):} se elige cómo enviar el aviso (carta documento, correo certificado, o un mensaje si la tienda lo acepta) y se sigue el envío con un número de rastreo.
    \item \textbf{Bloque 3 (nota, acordadas y plazos):} la tienda tiene fechas límite para responder y el reclamo cuenta con días hábiles e inhábiles.
    \item \textbf{Bloque 4 (demanda y primer proveído):} el reclamo debe decir quién es el reclamante, qué pasó y qué se pide; la tienda recibe una respuesta con instrucciones.
    \item \textbf{Bloque 5 (prueba y alegato):} se presentan facturas, fotos y testigos (prueba) y al final se hace un resumen de por qué se tiene razón (alegato).
\end{itemize}
\end{example}

% ============================================================
\chapter[Notificaciones: qué y por qué medio]{Notificaciones: qué se notifica y por qué medio}
% ============================================================

\section{Marco de la clase: la notificación dentro del proceso}

La clase anterior se ocupó de la demanda y de la contestación de demanda. Entre la demanda y su contestación se ubica la notificación; en rigor, la notificación está presente en todo el proceso. El tema se aborda con el fin de identificar las actividades que se realizan todos los días en la práctica.

La forma de notificación está regulada en el artículo 135 del Código Procesal.

\section{Artículo 135: resoluciones que se notifican por cédula}

Los supuestos del artículo 135 son originarios del Código, por lo que no han tenido mayores cambios.

\begin{formula}{Art.\ 135 CPCCN --- Notificación personal o por cédula}
Sólo serán notificadas personalmente o por cédula las siguientes resoluciones:
\begin{enumerate}[label=\arabic*)]
    \item La que dispone el traslado de la demanda, de la reconvención y de los documentos que se acompañen con sus contestaciones.
    \item La que dispone correr traslado de las excepciones y la que las resuelva.
    \item La que ordena la apertura a prueba y designa audiencia preliminar conforme al artículo 360.
    \item La que declare la cuestión de puro derecho, salvo que ello ocurra en la audiencia preliminar.
    \item Las que se dicten entre el llamamiento para la sentencia y ésta.
    \item Las que ordenan intimaciones, o apercibimientos no establecidos directamente por la ley, hacen saber medidas cautelares o su modificación o levantamiento, o disponen la reanudación de plazos suspendidos por tiempo indeterminado, o aplican correcciones disciplinarias.
    \item La providencia que hace saber la devolución del expediente, cuando no haya habido notificación de la resolución de alzada o cuando tenga por objeto reanudar plazos suspendidos por tiempo indeterminado.
    \item La primera providencia que se dicte después que un expediente haya vuelto del archivo de los tribunales, o haya estado paralizado o fuera de secretaría más de tres meses.
    \item Las que disponen vista de liquidaciones.
    \item La que ordena el traslado del pedido de levantamiento de embargo sin tercería.
    \item La que dispone la citación de personas extrañas al proceso.
    \item Las que se dicten como consecuencia de un acto procesal realizado antes de la oportunidad que la ley señala para su cumplimiento.
    \item Las sentencias definitivas y las interlocutorias con fuerza de tales y sus aclaratorias con excepción de las que resuelvan caducidad de la prueba por negligencia.
    \item La providencia que deniega los recursos extraordinarios.
    \item La providencia que hace saber el juez o tribunal que va a conocer en caso de recusación, excusación o admisión de la excepción de incompetencia.
    \item La que dispone el traslado del pedido de caducidad de la instancia.
    \item La que dispone el traslado de la prescripción en los supuestos del artículo 346, párrafos segundo y tercero.
    \item Las demás resoluciones de que se haga mención expresa en la ley o determine el Tribunal excepcionalmente, por resolución fundada.
\end{enumerate}
No se notificarán mediante cédula las decisiones dictadas en la audiencia preliminar a quienes se hallaren presentes o debieron encontrarse en ella.

Los funcionarios judiciales quedarán notificados el día de la recepción del expediente en su despacho. Deberán devolverlo dentro del tercer día, bajo apercibimiento de las medidas disciplinarias a que hubiere lugar.

No son aplicables las disposiciones contenidas en el párrafo precedente al Procurador General de la Nación, al Defensor General de la Nación, a los Procuradores Fiscales de la Corte Suprema, a los Procuradores Fiscales de Cámara, y a los Defensores Generales de Cámara, quienes serán notificados personalmente en su despacho.
\end{formula}

El inciso 1 dispone que el traslado de la demanda, de la reconvención y de los documentos acompañados con sus contestaciones se notifica por cédula: es decir, la demanda siempre va por cédula, salvo alguna excepción. De todos los supuestos del artículo, las cédulas de demanda y de sentencia son las dos que siempre van por cédula; existen varios supuestos básicos, pero estos dos son los principales.

\begin{important}
\textbf{Regla de oro:} la \textbf{demanda} y la \textbf{sentencia} siempre se notifican por \textbf{cédula} (salvo alguna excepción).
\end{important}

\section{Artículo 136: medios de notificación}

Las formas de notificación están en el artículo 136. Los medios allí enumerados son \textbf{taxativos}, en principio. Se trata de un artículo fundamental, cuya confusión suele traducirse en errores en los parciales.

\begin{formula}{Art.\ 136 CPCCN --- Medios de notificación}
En los casos en que este Código u otras leyes establezcan la notificación por cédula, ella también podrá realizarse por los siguientes medios:
\begin{enumerate}[label=\arabic*)]
    \item Acta notarial.
    \item Telegrama con copia certificada y aviso de entrega.
    \item Carta documento con aviso de entrega.
\end{enumerate}
La notificación de los traslados de demanda, reconvención, citación de personas extrañas al juicio, la sentencia definitiva y todas aquellas que deban efectuarse con entrega de copias, se efectuarán únicamente por cédula o acta notarial, sin perjuicio de la facultad reglamentaria concedida a la Corte Suprema de Justicia.

Se tendrá por cumplimentada la entrega de copias si se transcribe su contenido en la carta documento o telegrama.

La elección del medio de notificación se realizará por los letrados, sin necesidad de manifestación alguna en las actuaciones.

Los gastos que arrojen las notificaciones integrarán la condena en costas.

Ante el fracaso de una diligencia de notificación no será necesaria la reiteración de la solicitud del libramiento de una nueva, la que incluso podrá ser intentada por otra vía.
\end{formula}

\subsection{Medios utilizados en la práctica}

En la práctica los medios más conocidos y utilizados son la cédula, la carta documento y el telegrama. La norma menciona también los «medios de comunicación», que la docente señala no haber visto nunca aplicados; antes, en procesos colectivos, se usaban a veces edictos por radio.

\begin{table}[H]
\centering
\caption{Medios de notificación mencionados en clase}
\begin{tabularx}{\textwidth}{L{0.20\textwidth} Y Y}
\toprule
\textbf{Medio} & \textbf{Cuándo se usa} & \textbf{Observaciones} \\
\midrule
\textbf{Cédula (papel, electrónica o ley)} & Es el medio principal; la demanda y la sentencia siempre van por cédula. & Es la que más se conoce y la que se hace todos los días. \\
\addlinespace
\textbf{Carta documento} & Se puede mandar siempre, excepto sentencia o demanda y cuando hay copias que adjuntar. & Debe copiar la resolución completa y con datos precisos. \\
\addlinespace
\textbf{Telegrama} & Medio enumerado en el art.\ 136. & Rige la misma lógica que la carta documento. \\
\addlinespace
\textbf{Edicto} & Cuando se desconocen las partes o su domicilio. & Se publica en el boletín oficial de la jurisdicción. \\
\addlinespace
\textbf{Medios de comunicación} & Figuran en la norma, pero la docente no los vio aplicados. & Antes, edictos por radio en procesos colectivos. \\
\addlinespace
\textbf{Acta notarial} & No fue desarrollada en clase. & Está prevista en el art.\ 136. \\
\bottomrule
\end{tabularx}
\end{table}

\section{Edictos}

Los edictos se siguen utilizando. Corresponde notificar por edicto cuando se desconoce si existen más partes de las que se conocen o cuando no se sabe a quién ni adónde notificar: no es posible enviar una cédula porque no hay destinatario ni domicilio conocido.

\begin{example}{sucesión}
En una sucesión puede haber herederos que se desconocen; por esa razón corresponde la notificación por edicto.
\end{example}

El edicto se publica en el Boletín Oficial, en el boletín de la ciudad o en el de la provincia, según la jurisdicción de que se trate.

\begin{formula}{Art.\ 145 CPCCN --- Notificación por edictos (norma de referencia)}
Además de los casos determinados por este Código, procederá la notificación por edictos cuando se tratare de personas inciertas o cuyo domicilio se ignore. En este último caso, la parte deberá manifestar bajo juramento que ha realizado sin éxito las gestiones tendientes a conocer el domicilio de la persona a quien se deba notificar. Si resultare falsa la afirmación de la parte que dijo ignorar el domicilio, o que pudo conocerlo empleando la debida diligencia, se anulará a su costa todo lo actuado con posterioridad, y será condenada a pagar una multa de PESOS CINCUENTA (\$ 50) a PESOS QUINCE MIL (\$ 15.000).
\end{formula}

\section{Cuadro Cornell: notificaciones y medios}

\begin{longtable}{L{0.27\textwidth} L{0.67\textwidth}}
\toprule
\textbf{Preguntas / palabras clave} & \textbf{Notas / respuestas} \\
\midrule
\endhead

\textbf{¿Dónde está la notificación dentro del proceso?} &
Entre la demanda y la contestación, y en realidad en todo el proceso. Es lo que se hace todos los días, por eso la clase se centra en ella (art.\ 135 del Código). \\
\addlinespace

\textbf{¿Qué resoluciones siempre van por cédula (art.\ 135)?} &
La \textbf{demanda} (su traslado) y la \textbf{sentencia}, salvo alguna excepción. Son las dos resoluciones básicas que siempre se notifican por cédula. \\
\addlinespace

\textbf{¿Cuáles son los medios de notificación (art.\ 136) y por qué son taxativos?} &
Los enumerados en la ley, \textbf{taxativos} en principio: cédula, acta notarial, telegrama y carta documento. Los «medios de comunicación» casi no se ven en la práctica. El artículo resulta fundamental para el examen. \\
\addlinespace

\textbf{¿Cuándo corresponde un edicto y dónde se publica?} &
Cuando se desconoce si hay más partes o no se sabe a quién ni adónde notificar (por ejemplo, una sucesión). Se publica en el Boletín Oficial de la jurisdicción (nacional, de la ciudad o de la provincia). \\

\midrule
\multicolumn{2}{p{\dimexpr0.94\textwidth+2\tabcolsep\relax}}{\textbf{Resumen:} la notificación atraviesa todo el proceso. El artículo 135 enumera las resoluciones que van por cédula, y las de demanda y sentencia siempre se notifican de ese modo. El artículo 136 enumera los medios de notificación, que son taxativos en principio. Cuando se ignoran las partes o el domicilio, no hay a quién enviar una cédula y corresponde el edicto, publicado en el boletín de la jurisdicción.} \\
\bottomrule
\end{longtable}

% ============================================================
\chapter[La cédula y otros medios en la práctica]{La cédula y otros medios de notificación en la práctica}
% ============================================================

\section{Cédula ley y oficio ley 22.172: la «extraña jurisdicción»}

Existen distintas formas de cédula. Lo explicado no figura en los libros de modo diferente; se trata de preguntas frecuentes de examen.

\begin{definition}{Extraña jurisdicción}
Situación en la que el juzgado donde está radicado el juicio se encuentra en un lugar y el domicilio del demandado está en otro, es decir, fuera de la jurisdicción del tribunal. Cuando ello ocurre, la cédula es una \textbf{cédula ley}.
\end{definition}

\subsection{El caso de la Provincia de Buenos Aires}

Si se litiga, por ejemplo, en San Isidro y debe notificarse en La Matanza, \textbf{no} se trata de una cédula ley, porque la jurisdicción es la Provincia de Buenos Aires y no el departamento judicial.

\begin{example}{jurisdicción de la Provincia de Buenos Aires}
Mar del Plata es exactamente lo mismo que Quilmes: toda la Provincia de Buenos Aires es una sola jurisdicción. Chacabuco, San Isidro y Junín también están comprendidos en esa única jurisdicción.
\end{example}

\subsection{Oficio ley}

Un oficio dirigido, por ejemplo, a Avellaneda constituye un oficio ley. Es exactamente igual que un oficio común, salvo que debe llevar en la parte superior la inscripción \textbf{«Oficio ley 22.172»}.

Para la cédula ley el procedimiento exige que esté sellada por el juzgado, firmada por el secretario y por el juez, y que se consigne la competencia del juez. Además, debe estar firmada por el abogado, como todas las cédulas.

\begin{note}
Criterio de edición de las notas originales: en la grabación la docente dice «ley 22» y la transcripción automática la separa del resto de la cifra. Se completó como Ley 22.172 (comunicaciones entre tribunales de distintas jurisdicciones), norma a la que se refieren el «oficio ley» y la «cédula ley».
% TODO: contenido incompleto o ambiguo en las notas originales (número de ley completado por el editor de las notas)
\end{note}

\section{Trámite de la cédula en papel}

\subsection{Primero firmar y sellar}

La cédula en papel se firma y se lleva al juzgado, donde se sella.

\subsection{Dos opciones: gestoría u oficina de notificaciones}

Una vez sellada, existen dos opciones:

\begin{itemize}
    \item \textbf{Gestoría:} es el medio cómodo. Se sella en el juzgado, se deja en la gestoría y esta avisa cuando puede retirarse la cédula diligenciada.
    \item \textbf{Oficina de notificaciones:} quien tenga poco dinero o poco tiempo debe fijarse a qué oficina corresponde el domicilio donde va la cédula y llevarla directamente. Por ejemplo, si el domicilio de la persona está en Vicente López, corresponde la oficina de notificaciones de San Martín.
\end{itemize}

\subsection{Qué se gana yendo directamente}

Lo que se gana es \textbf{tiempo}, porque se evita todo lo que la gestoría hace para llegar a la oficina. Si la cédula se deja un día en la oficina, al otro día sale para el domicilio de la persona; luego hay que ir a retirarla, lo cual es fundamental para terminar el trámite.

Al dejar la cédula se entrega un \textbf{código de trámite}, con el cual se retira posteriormente. No debe perderse. Consiste en un papel blanco tamaño A4, de 4 por 4, que tiene un sello y un número. Un alumno señaló que en San Martín el código es electrónico: el código de barras.

\begin{important}
En el examen: la cédula en papel se sella en el juzgado y se lleva a la \textbf{oficina de notificaciones que corresponde} al domicilio; la gestoría es solo el medio cómodo. Conviene mencionar \textbf{ambas opciones}: afirmar únicamente que se lleva a la gestoría no resulta suficiente para aprobar.
\end{important}

La docente mencionó además que el PJN cambió.

\section{Cédula papel, cédula electrónica y seguimiento de su estado}

\subsection{Cuándo se usa la cédula en papel}

La cédula en papel se utiliza, en general, la primera vez, porque aún no hay abogado de la otra parte: el demandado tiene domicilio, no se presentó y no constituyó domicilio.

La diferencia con la cédula ley radica en que la cédula papel se dirige a un domicilio en la misma jurisdicción: si el demandado reside en Capital, el procedimiento es el descripto. La diligencia la realizan los abogados.

\subsection{Cómo consultar el estado de una cédula en papel}

\begin{enumerate}
    \item Ir al juzgado, a la mesa de entrada, y pedir \textbf{el libro de cédulas}, donde figuran todas las cédulas del juzgado entre una fecha y otra. Conviene ir cuando hayan pasado 10 o 15 días desde que se dejó la cédula.
    \item Identificar la cédula correspondiente, que tiene pegado un \textit{sticker} con un \textbf{código de barras}; se le saca una foto o se lo anota.
    \item Ingresar ese código en la página de la \textbf{Oficina de Mandamientos y Notificaciones de la Corte}, en el lugar indicado. La consulta es muy rápida.
\end{enumerate}

El sistema informa solo el resultado: «en trámite», «resultado positivo» o «resultado negativo». No indica, por ejemplo, que la persona no vive en el domicilio; eso se conocerá después. Aun así, es una buena forma de saber que la cédula ya salió.

\subsection{Cédula electrónica}

La cédula electrónica se utiliza cuando el demandado se presentó y tiene domicilio constituido (domicilio electrónico). Se confecciona y envía desde el sistema. La docente señaló que debía aprender el cambio introducido en el PJN.

\begin{table}[H]
\centering
\caption{Tipos de cédula}
\begin{tabularx}{\textwidth}{L{0.16\textwidth} Y Y Y}
\toprule
 & \textbf{Cédula ley} & \textbf{Cédula papel} & \textbf{Cédula electrónica} \\
\midrule
\textbf{¿Cuándo se usa?} & Cuando el domicilio está en otra jurisdicción («extraña jurisdicción»). & Primera notificación: el demandado aún no se presentó ni constituyó domicilio. & Cuando la otra parte se presentó y tiene domicilio electrónico constituido. \\
\addlinespace
\textbf{Requisitos y trámite} & Sello del juzgado, firma de secretario y juez, competencia del juez; firma del abogado; inscripción «Oficio ley 22.172». & Firma del abogado; sello del juzgado; oficina de notificaciones (o gestoría). & Se confecciona y envía desde el sistema. \\
\addlinespace
\textbf{Seguimiento} & Según la oficina de destino. & Libro de cédulas, código de barras y consulta en la página de la oficina de mandamientos y notificaciones de la Corte. & Queda registrada en el sistema. \\
\bottomrule
\end{tabularx}
\end{table}

\section{Familia y civil patrimonial: ¿se puede notificar por WhatsApp o por mail?}

Hay situaciones que en familia son mucho más permeables que en civil patrimonial. En un proceso patrimonial no se permite notificar por WhatsApp: hay que cumplir con el Código. En familia, en cambio, muchas veces se notifica por WhatsApp, por correo electrónico y de otras formas, incluso en los lugares de trabajo de las personas.

Todo lo que exceda el artículo 136 requiere \textbf{autorización del juzgado}; de lo contrario, la notificación es \textbf{nula}. No se puede enviar directamente un WhatsApp a la contraparte. Al pedir la autorización deben explicarse las razones por las cuales no es posible notificar de otra forma, y invocarse la celeridad procesal. En el práctico, por tratarse de patrocinios gratuitos, suele apelarse además a esa circunstancia, lo que resulta favorable en materia de notificaciones.

\begin{table}[H]
\centering
\caption{Notificar fuera del artículo 136}
\begin{tabularx}{\textwidth}{L{0.22\textwidth} Y Y}
\toprule
 & \textbf{Civil patrimonial} & \textbf{Familia} \\
\midrule
\textbf{WhatsApp / e-mail} & No se permiten: hay que cumplir con el Código. & Muchas veces se notifica así, con autorización. \\
\addlinespace
\textbf{Requisito} & Medios del art.\ 136, sin excepción. & Pedir autorización al juzgado; si no, la notificación es nula. \\
\addlinespace
\textbf{Argumentos útiles} & --- & Explicar por qué no se puede notificar por los medios comunes e invocar la celeridad procesal. \\
\bottomrule
\end{tabularx}
\end{table}

\section{Carta documento: usos, límites y contenido}

\subsection{Cuándo se puede mandar}

La carta documento se puede mandar siempre, excepto cuando se trata de \textbf{sentencia o demanda} y cuando hay \textbf{copias} que adjuntar, pues las copias no se pueden adjuntar. Si no hay copias, todas las demás resoluciones pueden notificarse por carta documento.

\subsection{Cómo se redacta}

El texto de la carta documento debe ser \textbf{autosuficiente}: debe copiarse toda la resolución, con su fecha y datos bien precisos. En la cédula, en cambio, muchas veces se abrevia diciendo «la siguiente resolución» y se adjunta la copia; en la carta documento eso no es posible.

\begin{note}
El art.\ 136 (transcripto en el capítulo anterior) prevé que se tiene por cumplida la entrega de copias si se transcribe su contenido en la carta documento o telegrama. Aun así, el traslado de demanda y la sentencia definitiva se notifican únicamente por cédula o acta notarial.
\end{note}

\section{Cuadro Cornell: la cédula y otros medios}

\begin{longtable}{L{0.27\textwidth} L{0.67\textwidth}}
\toprule
\textbf{Preguntas / palabras clave} & \textbf{Notas / respuestas} \\
\midrule
\endhead

\textbf{¿Qué es la cédula ley y el oficio ley 22.172?} &
Es la cédula (u oficio) que va a \textbf{extraña jurisdicción}. Dentro de la Provincia de Buenos Aires \textbf{no} es cédula ley, porque toda la provincia es una sola jurisdicción (Mar del Plata es lo mismo que Quilmes). Lleva la inscripción «Oficio ley 22.172», sello del juzgado, firmas de secretario, juez y abogado, y la competencia del juez. \\
\addlinespace

\textbf{¿Cómo se tramita una cédula en papel?} &
Se firma, se sella en el juzgado y se lleva a la \textbf{oficina de notificaciones} que corresponde al domicilio (o a la gestoría, que es el medio cómodo). Se gana tiempo yendo directamente a la oficina. Se retira con un código de trámite que no debe perderse. En el examen conviene nombrar ambas opciones. \\
\addlinespace

\textbf{¿Cuál es la diferencia entre cédula papel y cédula electrónica?} &
La cédula papel se usa en la primera notificación, cuando el demandado no se presentó ni constituyó domicilio. La electrónica se usa cuando la otra parte se presentó y tiene domicilio electrónico constituido, y se envía desde el sistema. \\
\addlinespace

\textbf{¿Cómo sigo el estado de una cédula?} &
Mesa de entrada, libro de cédulas, código de barras y consulta en la página de la oficina de mandamientos y notificaciones de la Corte. El resultado puede ser «en trámite», «positivo» o «negativo». \\
\addlinespace

\textbf{¿Puedo notificar por WhatsApp o mail?} &
En civil patrimonial \textbf{no}: hay que cumplir el art.\ 136. En familia es más permeable, pero todo lo que exceda el 136 requiere \textbf{autorización del juzgado}; si no, la notificación es nula. Se explican las razones y se invoca la celeridad procesal. \\
\addlinespace

\textbf{¿Cuándo sirve la carta documento y qué debe contener?} &
Siempre, salvo para \textbf{demanda o sentencia} y cuando hay \textbf{copias} que adjuntar. Debe ser autosuficiente: copiar toda la resolución, con fecha y datos precisos. \\

\midrule
\multicolumn{2}{p{\dimexpr0.94\textwidth+2\tabcolsep\relax}}{\textbf{Resumen:} la cédula puede ser ley (extraña jurisdicción), papel (primera notificación, sin domicilio constituido) o electrónica (domicilio constituido). La cédula en papel se sella en el juzgado y se lleva a la oficina de notificaciones; su estado se sigue con el libro de cédulas y el código de barras. Lo que exceda el art.\ 136, como WhatsApp o mail, exige autorización del juzgado, con mayor flexibilidad en familia. La carta documento no sirve para demanda ni sentencia ni cuando hay copias, y su texto debe bastarse a sí mismo.} \\
\bottomrule
\end{longtable}

% ============================================================
\chapter[Notificación por nota, acordadas y plazos]{Notificación por nota, acordadas y plazos}
% ============================================================

\section{Notificación por nota (ministerio de la ley): martes y viernes}

Los días de nota son \textbf{martes y viernes}. En el amparo, en cambio, son casi todos los días. La ley de amparo que regula ese régimen no se trabaja en la clase y debe leerse por cuenta propia.

\begin{formula}{Art.\ 133 CPCCN --- Principio general (norma de referencia)}
Salvo los casos en que procede la notificación por cédula y sin perjuicio de lo dispuesto en el artículo siguiente, las resoluciones judiciales quedarán notificadas en todas las instancias los días martes y viernes. Si uno de ellos fuere feriado, la notificación tendrá lugar el siguiente día de nota.

No se considerará cumplida tal notificación:
\begin{enumerate}[label=\arabic*)]
    \item Si el expediente no se encontrare en el tribunal.
    \item Si hallándose en él, no se exhibiere a quien lo solicita y se hiciera constar tal circunstancia en el libro de asistencia por las personas indicadas en el artículo siguiente, que deberá llevarse a ese efecto.
\end{enumerate}
Incurrirá en falta grave el prosecretario administrativo que no mantenga a disposición de los litigantes o profesionales el libro mencionado.
\end{formula}

\section{Acordadas: qué son y cómo afectan los plazos}

\begin{definition}{Acordada}
Decisión adoptada por los ministros de la Corte, que regula la actividad jurisdiccional. Por ejemplo, establece que un día es feriado o inhábil, que desde cierta fecha rige la feria, o que un juzgado está cerrado por reforma. Funciona como una ley que regula a la Corte y a todos los operadores.
\end{definition}

Si falta regular algún aspecto judicial, es probable que se haga por acordada: mudanzas de juzgado, todo lo relativo a la organización interna y los días inhábiles.

\subsection{Acordadas y caídas del sistema}

Hace unos años el PJN se caía con frecuencia, quedando sin funcionar durante dos o tres horas; en la actualidad funciona mucho mejor. En aquel entonces, cada tanto una acordada declaraba un día inhábil.

Que un día sea inhábil significa que deben \textbf{contarse de nuevo los plazos}.

\begin{example}{recómputo por día inhábil}
Si un plazo debía presentarse el lunes (contando martes, jueves, viernes y el lunes), y se declara inhábil el lunes, el vencimiento se corre al martes. Lo mismo ocurre durante la semana con los feriados.
\end{example}

\begin{important}
Antes de dar un plazo por vencido, \textbf{debe revisarse si hay acordadas}: un día declarado inhábil obliga a \textbf{recontar} el plazo. Cuando el plazo está justo y surge la duda, la revisión de las acordadas puede resultar decisiva.
\end{important}

\section{Dejar nota, expedientes «en letra» y cómputo de plazos}

\subsection{Del papel a la virtualidad}

El régimen de la nota fue concebido para la forma analógica y hoy se adapta a la virtual. Antes, con todo en papel, había que ir a los juzgados los martes y viernes; se entregaba un cuaderno y allí se dejaba constancia de que se había concurrido y de que el expediente no se encontraba. Eso era \textbf{dejar nota}.

Era fundamental porque los expedientes en papel se movían: podían estar en la defensoría, en la fiscalía o en despacho, y no se sabía qué había ocurrido. Hoy todo, o casi todo, está en la grilla.

\begin{definition}{Dejar nota}
Constancia de que se miró el expediente en el día de nota. Hoy se realiza de manera virtual, desde el usuario del abogado: hay un botón superior, «Dejar nota», que trae todos los expedientes habilitados para ello. Los alumnos no pueden dejar nota desde su usuario.
\end{definition}

\begin{definition}{Expediente «en letra»}
Expediente que se puede consultar. Antes, «en letra» indicaba que lo podía ver mesa de entrada, y cuando estaba en despacho se indicaba esa circunstancia. En el PJN, los expedientes que figuran «en letra» no permiten dejar nota.
\end{definition}

Respecto de la digitalización, los expedientes de 2017 hacia atrás suelen estar en papel. Depende del juzgado: algunos lo digitalizaron todo y son los que mejor trabajan; otros digitalizaron de a poco, y otros no digitalizaron. Hay causas muy voluminosas: en daños y perjuicios, expedientes de siete cuerpos; en comerciales y quiebras, de 20 o 30 cuerpos.

\subsection{¿Cuándo me entero? El cómputo del plazo}

Si la resolución sale un lunes a la tarde, la cuestión es cuándo se toma nota y desde cuándo se cuenta el plazo.

En el esquema analógico, el lunes no se veía la resolución; el martes por la mañana se iba al juzgado y se tomaba conocimiento allí. El primer día del plazo era el miércoles, y los plazos se contaban seguidos. Era complejo porque se empezaba a contar, el expediente desaparecía y había que volver a contar.

Con el sistema actual, si el proveído sale el martes, queda notificado el viernes siguiente si se dejó nota. Si se dejó nota a las 10:00 y el proveído sale a las 14:30, se está cubierto, y el plazo empieza a correr el viernes. Existe una ventaja: se puede trabajar el caso martes, miércoles, jueves y viernes, ganando tres días para avisar al cliente, reunirse y obtener las firmas necesarias, porque la resolución ya se ve en la grilla.

Si la resolución sale un viernes, se cuenta de modo análogo: el lunes es el primer día. Si no se deja nota, el plazo cuenta desde el miércoles.
% TODO: contenido incompleto o ambiguo en las notas originales (cómputo cuando no se deja nota)

\begin{table}[H]
\centering
\caption{Ejemplos de cómputo citados por la docente}
\begin{tabularx}{\textwidth}{L{0.28\textwidth} Y L{0.18\textwidth}}
\toprule
\textbf{Situación} & \textbf{Qué ocurre} & \textbf{Primer día de plazo} \\
\midrule
\textbf{Expediente en papel: la resolución sale el lunes} & El martes a la mañana se iba al juzgado y se notificaban allí. & El miércoles. \\
\addlinespace
\textbf{Con nota: se deja nota a las 10:00 y el proveído sale a las 14:30} & «Se está cubierto»: se puede ir trabajando el caso desde el martes. & El viernes (se ganan tres días para trabajar). \\
\addlinespace
\textbf{La resolución sale un viernes} & Se cuenta de modo análogo. & El lunes. \\
\addlinespace
\textbf{No se dejó nota} & Según la docente, sin nota el plazo cuenta desde el miércoles. & El miércoles. \\
\bottomrule
\end{tabularx}
\end{table}

\begin{note}
La propia docente reconoce que este es el punto más débil de su explicación y pidió simularlo en un expediente real: tomar un expediente, buscar con el usuario el botón de dejar nota, observar cuándo sale la resolución y simular el cómputo, para plantear las dudas en la clase siguiente. La regla de fondo es la del art.\ 133 (transcripto más arriba): las resoluciones quedan notificadas los martes y viernes (o el siguiente día de nota si uno es feriado), salvo los casos de cédula.
\end{note}

\subsection{Plazos del amparo}

En el amparo los días de nota son casi todos los días, excepto los lunes en que se deja nota. Los plazos son muy cortos: de 24 horas, y hasta de 12 horas; son horas, no días. Las apelaciones son de 48 horas. Quien tome un amparo sin experiencia debe dar suma importancia a los plazos y leer la ley.

Criterio práctico: \textit{ante la duda, contar corto}. Si hay dudas sobre si un plazo vence un lunes o un martes, debe tomarse el lunes.

\subsection{¿Siempre hay que dejar nota?}

En la práctica del grupo no se le presta mucha atención a dejar nota, porque siempre hay alguien del equipo controlando las resoluciones y existe un sistema automatizado de control. Sin embargo, es necesario hacerlo, porque puede pasar algo por alto.

\begin{important}
\textbf{Dejar nota} es la constancia de que se miró el expediente en el día de nota. Si el expediente está «en letra», el sistema no permite dejarla. \textbf{Ante la duda, contar corto.}
\end{important}

\section{Cuadro Cornell: nota, acordadas y plazos}

\begin{longtable}{L{0.27\textwidth} L{0.67\textwidth}}
\toprule
\textbf{Preguntas / palabras clave} & \textbf{Notas / respuestas} \\
\midrule
\endhead

\textbf{¿Qué días se notifica por nota?} &
Martes y viernes (o el siguiente día de nota si uno es feriado). En el amparo, casi todos los días. \\
\addlinespace

\textbf{¿Qué es una acordada y por qué importa?} &
Es una decisión de la Corte que regula la actividad jurisdiccional (feriados, ferias, días inhábiles, organización interna). Un día inhábil obliga a \textbf{recontar los plazos}; revisar las acordadas «puede salvar la vida». \\
\addlinespace

\textbf{¿Qué es «dejar nota» y qué significa «en letra»?} &
Dejar nota es la constancia de que se consultó el expediente en el día de nota. En el PJN hay un botón «Dejar nota»; si el expediente está \textbf{«en letra»} no permite dejarla. \\
\addlinespace

\textbf{¿Desde cuándo corre el plazo si dejo nota?} &
Si se deja nota y el proveído sale el martes, queda notificado el viernes y desde allí cuenta el plazo, ganando tres días para trabajar el caso. Si sale un viernes, el lunes es el primer día. Sin nota, según la docente, cuenta desde el miércoles. La propia docente señala que es un punto débil y pide simularlo en un expediente real. \\
\addlinespace

\textbf{¿Qué plazos tiene el amparo?} &
Muy cortos: de 24 y hasta 12 horas; las apelaciones, 48 horas. Hay que leer la ley de amparo. \\
\addlinespace

\textbf{¿Qué criterio se aplica ante la duda sobre un plazo?} &
Contar corto: si hay duda sobre si vence un lunes o un martes, se toma el lunes. \\

\midrule
\multicolumn{2}{p{\dimexpr0.94\textwidth+2\tabcolsep\relax}}{\textbf{Resumen:} las resoluciones no notificadas por cédula quedan notificadas por nota los martes y viernes. Las acordadas de la Corte pueden declarar días inhábiles y obligan a recontar plazos. Dejar nota es la constancia de haber mirado el expediente en día de nota y se hace con un botón del PJN, salvo que el expediente esté «en letra». El cómputo del plazo es un punto delicado que conviene simular en un expediente real; en el amparo los plazos son de horas. Ante la duda, siempre se cuenta corto.} \\
\bottomrule
\end{longtable}

% ============================================================
\chapter[Demanda y primer proveído]{Demanda y primer proveído}
% ============================================================

\section{La demanda (art.\ 330) y cómo se inicia el expediente}

\subsection{Cierre del tema notificaciones y tarea}

Se dieron por vistas las formas de notificación, tratadas con ejemplos muy aleatorios. Más adelante, cuando se redacte la demanda, se volverá sobre el tema. La docente advirtió que la nota no había quedado clara y pidió leerla, simular en un expediente el procedimiento de tomar nota (con el botón de dejar nota del usuario), anotar las dudas y leer el Código, donde está explicada.

\subsection{Cómo se arma una demanda}

El artículo 330 indica cómo hacer una demanda, qué poner y qué no poner. Ante la duda se sigue el orden: \textbf{encabezado, derecho, objeto, competencia}. Funciona como un esquema para armar. Por regla general, \textbf{todo lo que no esté en la demanda no se puede agregar}.

\begin{formula}{Art.\ 330 CPCCN --- Forma de la demanda (contenido esencial)}
La demanda será deducida por escrito y contendrá: nombre y domicilio del demandante; nombre y domicilio del demandado; la cosa demandada, designándola con toda exactitud; los hechos en que se funde, explicados claramente; el derecho expuesto sucintamente, evitando repeticiones innecesarias; y la petición en términos claros y positivos.
\end{formula}

\begin{note}
El artículo no fue leído en clase y no figuraba en las fuentes consultadas para las notas; se consignó su contenido esencial. Debe verificarse el texto íntegro en el Código vigente.
% TODO: contenido incompleto o ambiguo en las notas originales (texto íntegro del art. 330 no verificado)
\end{note}

\subsection{Interposición de la demanda}

Para interponer la demanda se envía un mail a la Cámara Civil, que informa el número de juzgado y de expediente. Se sube la demanda y la documental. La demanda debe contar con la documental como prueba, además de todas las otras pruebas que se pretendan producir. En el primer proveído, el juzgado ordena: «notifíquese».

\section[Primer proveído y forma de notificar]{Primer proveído y forma de notificar}

\subsection{La palabra del proveído indica cómo notificar}

Cuando un proveído dice \textbf{«notifíquese»}, corresponde \textbf{cédula}, sin duda. Si dice \textbf{«hágase saber»}, corresponde, en general, notificar por \textbf{nota}. Si dice \textbf{«traslado»}, o «notifíquese traslado», el traslado va por nota; recuérdese que el traslado de la demanda siempre va por cédula.

Algunos proveídos no dicen nada. En ese caso no se puede notificar: si el proveído no lo dice, se necesita autorización del juez para hacer una cédula o cualquier otra notificación, y en principio puede pedirse explicación.

\subsection{Cuando el proveído «no dice nada»}

Por regla general, comprobada en la práctica, cuando el proveído no contiene ninguna de esas tres expresiones se está \textbf{rechazando} algo.

\begin{example}{proveído sin las tres palabras}
«Atento a lo solicitado, no ha lugar al oficio. No se libre el oficio dispuesto a FJA 13.» El proveído no contiene ninguna de las tres expresiones, y ante él puede plantearse una revocatoria.
\end{example}

Respecto de desde cuándo se cuenta un proveído que no contiene ninguna de las tres palabras, se lo cuenta como nota, desde que sale. La docente lo cuenta desde el día en que lo leyó: el día siguiente es el primero. Esto es especialmente relevante si hay otra parte, porque el juzgado en general lo dejará pasar, pero la contraparte puede aprovechar la confusión.

\begin{table}[H]
\centering
\caption{Qué palabra figura en el proveído y cómo se notifica}
\begin{tabularx}{\textwidth}{L{0.25\textwidth} L{0.22\textwidth} Y}
\toprule
\textbf{El proveído dice\ldots} & \textbf{Se notifica por\ldots} & \textbf{Observaciones} \\
\midrule
\textbf{«Notifíquese»} & Cédula & «No hay duda.» \\
\addlinespace
\textbf{«Traslado» / «notifíquese traslado»} & Nota & El traslado de la demanda siempre va por cédula. \\
\addlinespace
\textbf{«Hágase saber»} & Nota, en general & --- \\
\addlinespace
\textbf{No dice nada} & No se puede notificar sin autorización del juez & Suele ser un rechazo; se puede plantear revocatoria y conviene contar el plazo desde que sale. \\
\bottomrule
\end{tabularx}
\end{table}

\subsection{Documentar la prueba y cédula de la demanda}

Una vez dictado el proveído, se confecciona la cédula correspondiente, porque la demanda se notifica por cédula, y la cédula sale.

\section{Ampliación de demanda y de prueba; hechos nuevos}

\subsection{El problema: omisiones posteriores}

\begin{example}{escena planteada por la docente}
Al releer la prueba en el estudio, el abogado advierte que olvidó algo fundamental. O bien el cliente dice que olvidó contar cierta parte, o muestra un papel que encontró en el cajón de su madre. Era lo que debía presentarse el primer día.
\end{example}

Ante esa situación se puede pedir \textbf{ampliación de demanda} y \textbf{ampliación de prueba}, pero con mucha rapidez.

\subsection{¿Hasta cuándo? Depende de si la cédula ya salió}

\begin{itemize}
    \item \textbf{Si la cédula ya se notificó} (volvió al juzgado con la leyenda de notificación), no se puede ampliar. Entre que la cédula sale del juzgado y llega, hay un margen para presentar el escrito, y el juzgado volverá a ordenar el traslado.
    \item \textbf{Si la cédula no salió,} no hay problema: se presenta el escrito y sale todo junto en la cédula.
    \item \textbf{Si la cédula salió pero aún no se notificó,} hay que actuar antes de que vuelva. Si la segunda cédula llega después a la otra parte, esta puede sostener que la demanda está notificada con el texto original, de modo que la ampliación no se toma.
\end{itemize}

El tiempo disponible para modificar es el del primer proveído, mientras la cédula se encuentra en trámite; una vez que se fue y volvió notificada, ya no se puede.

\subsection{Hechos nuevos en la contestación}

Al contestar el traslado de la contestación de demanda, se puede ampliar la prueba si hay un \textbf{hecho nuevo}. Lo que ya se dijo en la demanda y lo que se quiso probar, ya está dicho. Si, contestada la demanda, la parte advierte un hecho nuevo, puede manifestarlo («esto es un hecho nuevo, no estaba en mi demanda») y adjuntar documental. Debe tratarse de algo nuevo, que sorprendió o que no se conocía, y no de algo de la demanda original.

Existen muchas maniobras para hacer pasar un hecho viejo por hecho nuevo. Si la contraparte tiene un abogado que domine el derecho procesal, las dificultará, porque el Código es muy claro. A veces los clientes no conocen los papeles que tienen, y aparecen después.

\begin{formula}{Art.\ 365 CPCCN --- Hechos nuevos (norma de referencia)}
Cuando con posterioridad a la contestación de la demanda o reconvención, ocurriese o llegase a conocimiento de las partes algún hecho que tuviese relación con la cuestión que se ventila, podrán alegarlo hasta cinco días después de notificada la audiencia prevista en el artículo 360 del presente Código, acompañando la prueba documental y ofreciendo las demás de las que intenten valerse. Del escrito en que se alegue, si lo considerare pertinente, se dará traslado a la otra parte, quien, dentro del plazo para contestarlo, podrá también alegar otros hechos en contraposición a los nuevos alegados. El juez decidirá en la audiencia del artículo 360 la admisión o el rechazo de los hechos nuevos.
\end{formula}

\begin{formula}{Art.\ 334 CPCCN --- Hechos no invocados en la demanda o contrademanda (norma de referencia)}
Cuando en el responde de la demanda o de la reconvención se alegaren hechos no invocados en la demanda o contrademanda, los demandantes o reconvinientes según el caso podrán ofrecer prueba y agregar la documental referente a esos hechos, dentro de los 5 días de notificada la providencia respectiva. En tales casos se dará traslado de los documentos a la otra parte, quien deberá cumplir la carga que prevé el art.\ 356 inc.\ 1).
\end{formula}

\begin{important}
Ampliar demanda o prueba es posible \textbf{mientras la cédula del traslado no se haya diligenciado}. Una vez notificada, solo cabe ofrecer prueba sobre \textbf{hechos realmente nuevos} (por ejemplo, los alegados al contestar).
\end{important}

\section{Cuadro Cornell: demanda y primer proveído}

\begin{longtable}{L{0.27\textwidth} L{0.67\textwidth}}
\toprule
\textbf{Preguntas / palabras clave} & \textbf{Notas / respuestas} \\
\midrule
\endhead

\textbf{¿Qué contiene la demanda (art.\ 330)?} &
Ante la duda se sigue el orden: encabezado, derecho, objeto, competencia (y, según el contenido esencial consignado, hechos y petición), más la documental y las demás pruebas. Lo que no está en la demanda, por regla general, no se puede agregar. \\
\addlinespace

\textbf{¿Cómo se inicia el expediente?} &
Se envía un mail a la Cámara Civil, que informa el número de juzgado y de expediente; se sube la demanda con la documental y las demás pruebas. En el primer proveído el juzgado ordena «notifíquese». \\
\addlinespace

\textbf{¿Cómo notifico según lo que diga el proveído?} &
«Notifíquese» = cédula. «Hágase saber» = nota, en general. «Traslado» = nota (salvo el de la demanda, que va por cédula). Si no dice nada, no se puede notificar sin autorización del juez; suele ser un rechazo y puede plantearse revocatoria. \\
\addlinespace

\textbf{¿Hasta cuándo puedo ampliar la demanda o la prueba?} &
\textbf{Rapidísimo} y solo mientras la cédula del traslado no se haya notificado. Si no salió, se presenta el escrito y sale todo junto; si ya se notificó, no se puede. \\
\addlinespace

\textbf{¿Qué es un hecho nuevo y cuándo se puede alegar?} &
Un hecho que sorprendió o que la parte no conocía, distinto de lo ya dicho en la demanda. Al contestarse la demanda se puede alegarlo y adjuntar documental. Hay maniobras para hacer pasar un hecho viejo por nuevo, pero el Código es muy claro. \\

\midrule
\multicolumn{2}{p{\dimexpr0.94\textwidth+2\tabcolsep\relax}}{\textbf{Resumen:} la demanda se arma siguiendo el art.\ 330 y no puede ampliarse con lo omitido salvo en el margen previo a que se notifique la cédula. El primer proveído indica cómo notificar según su fórmula («notifíquese», «hágase saber», «traslado»); si no dice nada, suele ser un rechazo. Pasada la notificación, solo cabe ofrecer prueba sobre hechos realmente nuevos.} \\
\bottomrule
\end{longtable}

% ============================================================
\chapter[Prueba y alegatos]{Prueba y alegatos}
% ============================================================

\section{Ofrecimiento y fundamentación de la prueba}

\subsection{La prueba es la que enumera el Código}

La prueba que corresponde a la demanda está detallada taxativamente en el Código. Si se pretende producir otras pruebas distintas de las allí enunciadas, hay que pedir autorización.

\subsection{Explicar para qué se ofrece cada prueba}

Todas las pruebas deben estar explicadas: debe decirse por qué se ofrecen.

\begin{itemize}
    \item En la \textbf{prueba informativa}, hay que aclarar qué debe contestar el informe; no basta con poner «oficiar».
    \item En la \textbf{prueba pericial}, hay que proponer \textbf{los puntos de pericia}.
    \item En la \textbf{prueba testimonial}, hay que indicar para qué se cita a cada testigo.
\end{itemize}

\begin{example}{cómo fundamentar cada prueba}
Si se citan tres testigos, debe decirse para qué: para que den detalles sobre la dinámica familiar. No hace falta decir mucho más, pero tiene que estar. Si se pide un oficio, se afirma que informe cuánto cobra determinada persona. Si se pide una pericia psiquiátrica, es para ver si la persona está bien de la cabeza o no.
\end{example}

Los puntos de pericia, cuando el juez ordene la prueba, se copian y pegan en el oficio, que se lleva al perito o a la entidad correspondiente. Como dato útil, Mercado Pago tiene dirección de e-mail.

\begin{formula}{Art.\ 333 CPCCN --- Agregación de la prueba documental y ofrecimiento de la confesional (norma de referencia)}
Con la demanda, reconvención y contestación de ambas, deberá acompañarse la prueba documental y ofrecerse todas las demás pruebas de que las partes intentaren valerse. Cuando la prueba documental no estuviere a su disposición, la parte interesada deberá individualizarla, indicando su contenido, el lugar, archivo, oficina pública o persona en cuyo poder se encuentra. Si se tratare de prueba documental oportunamente ofrecida, los letrados patrocinantes, una vez interpuesta la demanda, podrán requerir directamente a entidades privadas, sin necesidad de previa petición judicial, y mediante oficio en el que se transcribirá este artículo, el envío de la pertinente documentación o de su copia auténtica, la que deberá ser remitida directamente a la secretaría, con transcripción o copia del oficio. Si se ofreciera prueba testimonial se indicará qué extremos quieren probarse con la declaración de cada testigo. Tratándose de prueba pericial la parte interesada propondrá los puntos de pericia.
\end{formula}

\section{Apertura a prueba y audiencia 360: los 40 días}

\subsection{El proveído de apertura a prueba}

Cumplidas las etapas de mediación, demanda y contestación de demanda, y superados los hechos nuevos, el juez, a pedido de parte, \textbf{abre a prueba} cuando recibe la demanda contestada. La apertura se articula mediante la \textbf{audiencia 360}.

Se necesita un proveído que ordene la apertura de la prueba, antes, durante o después de la celebración de la audiencia. Se requiere un proveído con \textbf{fecha cierta} para contar los \textbf{40 días}: desde la audiencia corren 40 días hábiles, durante los cuales se debe producir toda la prueba de la demanda y controlar que el demandado cumpla con la suya.

\begin{formula}{Art.\ 360 CPCCN --- Audiencia preliminar}
A los fines del artículo precedente el juez citará a las partes a una audiencia, que presidirá, con carácter indelegable. Si el juez no se hallare presente no se realizará la audiencia, debiéndose dejar constancia en el libro de asistencia. En tal acto:
\begin{enumerate}[label=\arabic*)]
    \item Invitará a las partes a una conciliación o a encontrar otra forma de solución de conflictos.
    \item Recibirá las manifestaciones de las partes con referencia a lo prescripto en el artículo 361 del presente Código, debiendo resolver en el mismo acto.
    \item Oídas las partes, fijará los hechos articulados que sean conducentes a la decisión del juicio sobre los cuales versará la prueba.
    \item Recibirá la prueba confesional si ésta hubiera sido ofrecida por las partes. La ausencia de uno de todos los absolventes, no impedirá la celebración de la audiencia preliminar.
    \item Proveerá en dicha audiencia las pruebas que considere admisibles y concentrará en una sola audiencia la prueba testimonial, la que se celebrará con presencia del juez en las condiciones establecidas en este capítulo. Esta obligación únicamente podrá delegarse en el secretario o en su caso, en el prosecretario letrado.
    \item Si correspondiere, decidirá en el acto de la audiencia que la cuestión debe ser resuelta como de puro derecho con lo que la causa quedará concluida para definitiva.
\end{enumerate}
\end{formula}

\begin{formula}{Art.\ 367 CPCCN --- Plazo de producción de prueba (norma de referencia)}
El plazo de producción de prueba será fijado por el juez, y no excederá de cuarenta días. Dicho plazo es común y comenzará a correr a partir de la fecha de celebración de la audiencia prevista en el artículo 360 del presente Código.
\end{formula}

\subsection{Caducidad y negligencia}

Si el demandado no cumple con una prueba, se puede pedir la \textbf{caducidad} o la \textbf{negligencia}. La docente indicó que existe un buen cuadro para diferenciar la negligencia, pero que para el examen basta con leer del Código en qué consiste cada una. Hay caducidad y negligencia respecto de las pruebas de la otra parte, que pueden pedirse, y la otra parte puede pedirlas a su vez.

\subsection{Cómo es en la práctica}

En general, sobre todo en familia, se convoca a una audiencia de conciliación antes de la audiencia. Si no hay conciliación, el juez o el secretario labra un acta, del tipo «vinieron las partes, no hubo acuerdo». Muy pocos juzgados incluyen en esa acta la apertura a prueba, lo cual sería lo más cómodo. En un caso mencionado por los alumnos, la apertura a prueba se dispuso en la propia audiencia.

Si en la audiencia se informa que no hay acuerdo, al terminar hay que presentar un escrito pidiendo que se abra a prueba, y el juzgado dicta el llamado \textbf{proveído de prueba}.

\begin{note}
Fragmento parcialmente inaudible (42:27 a 43:22 de la grabación): los alumnos comentan un caso del práctico en el que la clienta manifestó que no quería continuar con el juicio, se evalúa si corresponde la renuncia y se aclara que la apertura a prueba se dispuso «por el tema de responsabilidad». La grabación no permite reconstruir el diálogo textualmente, por lo que se consignan solo los datos inteligibles.
% TODO: contenido incompleto o ambiguo en las notas originales
\end{note}

\section{Proveído de prueba, control y apelación: caducidad, negligencia y prueba denegada}

\subsection{Qué es el proveído de prueba}

El secretario toma la demanda y la contestación de demanda, verifica qué prueba ofreció cada parte y las provee: por ejemplo, «Al pedido de oficio, líbrese. A la testimonial, fíjase fecha para el día\ldots». Provee toda la prueba de la demanda y de la contestación.

\subsection{Controlar y reclamar}

Muchas veces se omiten pruebas o se las pasa por alto, lo cual implica más trabajo para el abogado. Con el proveído de prueba hay que tomar el escrito propio y controlar qué se proveyó y qué no; lo que no se dijo, debe pedirse que se diga.

Si el juzgado rechaza una prueba (por ejemplo, «no ha lugar al oficio a Mercado Pago»), no se puede interponer nada, porque es una \textbf{facultad ordenatoria del juez}, y las resoluciones dictadas en ejercicio de ella son inapelables, aunque la prueba fuera importantísima para la parte. Se puede pedir que pase a la Cámara.

Cuando el caso llega a la sentencia y se interpone el recurso de apelación, es el momento de plantear lo acumulado desde la audiencia 360: en los agravios se plantea por qué se falló en contra, por la prueba que no se permitió producir, entre otros puntos del recurso, que no será el único.

\subsection{Negligencia y caducidad distintas de la prueba denegada}

\textbf{La negligencia} se da cuando una de las partes no produjo la prueba y dejó vencer los 40 días. La otra parte, o el juzgado, declara la caducidad o la negligencia, para esa parte o para la contraria. \textbf{Eso se puede apelar}, porque es una resolución.

En cambio, en el proveído de prueba el juzgado indica que no concederá una prueba: eso es facultativo, ordenatorio. Igualmente, el planteo puede reiterarse en los alegatos. El proveído debe controlarse, y el reclamo por la prueba no producida se reserva \textbf{hasta la sentencia}.

\begin{table}[H]
\centering
\caption{Negligencia o caducidad de la prueba frente a prueba denegada en el proveído}
\begin{tabularx}{\textwidth}{L{0.17\textwidth} Y Y}
\toprule
 & \textbf{Negligencia o caducidad de la prueba} & \textbf{Prueba denegada en el proveído de prueba} \\
\midrule
\textbf{¿Qué es?} & La parte no produjo su prueba y dejó vencer el plazo de 40 días. & El juzgado no provee o rechaza una prueba ofrecida («no ha lugar al oficio»). \\
\addlinespace
\textbf{Naturaleza} & Resolución del juzgado sobre el cumplimiento de la carga probatoria. & Facultad ordenatoria del juez. \\
\addlinespace
\textbf{¿Es apelable?} & Sí, porque es una resolución. & No es apelable en ese momento; se reclama en la apelación de la sentencia y, en lo pertinente, en los alegatos. \\
\addlinespace
\textbf{¿Qué hace el abogado?} & Pedirla contra la otra parte o defenderse si se la piden. & Controlar el proveído, pedir lo que falte y guardar el reclamo hasta la sentencia. \\
\bottomrule
\end{tabularx}
\end{table}

\section{Alegatos: estructura, plazos y visibilidad}

\subsection{Cuándo no se hacen alegatos}

Producida la prueba, sin caducidad ni negligencia y cumplidos los 40 días, viene la etapa de los \textbf{alegatos}.

En familia, en general, no se hacen alegatos, por tratarse de un proceso especial. Tampoco en otros procesos especiales, como el amparo, porque se entiende que retrasan el proceso y carecen de mayor sentido. En alimentos, por ejemplo, no hay tanta prueba (paga o no paga, trabaja o no trabaja). Si se trata de una cuestión muy especial, el juez puede ordenar que se produzcan los alegatos.

\subsection{Qué es un alegato y por qué importa}

\begin{definition}{Alegato}
Resumen que cada parte hace, con fuerza y concentración, de lo que fue el juicio: lo que probó con sus pruebas, lo que intentó probar y no pudo porque no se lo permitieron, y lo que la otra parte no probó.
\end{definition}

En general, para la sentencia es lo único que se lee, y más aún si el caso va a Cámara: allí se lee el alegato y la sentencia y se resuelve. Por eso es una de las partes más importantes del proceso. La estructura es la siguiente:

\begin{enumerate}
    \item \textbf{¿Qué probé yo?}
    \item \textbf{¿Qué no pude probar?}
    \item \textbf{¿Qué intentó probar la otra parte y tampoco probó?}
\end{enumerate}

No conviene esconder lo que le salió bien a la otra parte, porque negarlo resulta extraño, aunque siempre es bueno restarle un poco de valor. Lo más importante es \textbf{cómo afectó} la prueba que produjo la otra parte.

\subsection{Plazo: seis días por parte}

El juez dicta una resolución para que las partes pasen a alegar. El secretario controla que esté cumplida la prueba de ambas partes y dispone: «Pasen los autos a alegar». El plazo es de \textbf{seis días por parte}, no por separado. Si hay tres partes, a la última le tocan 18 días. Alegan primero las actoras y después las demandadas.

Hoy este esquema pierde sentido, porque el expediente está en la grilla y los seis días corren para todos. Antes, el abogado se llevaba el expediente al estudio, lo devolvía a los seis días, lo llevaba el otro y lo devolvía con su alegato a los seis días.

\begin{formula}{Art.\ 482 CPCCN --- Agregación de las pruebas. Alegatos (norma de referencia)}
Producida la prueba, el prosecretario administrativo, sin necesidad de gestión alguna de los interesados, o sin sustanciarla si se hiciera, ordenará que se agregue al expediente. Cumplido este trámite el prosecretario administrativo pondrá los autos en secretaría para alegar; esta providencia se notificará por cédula y una vez firme se entregará el expediente a los letrados por su orden y por el plazo de seis días a cada uno, sin necesidad de petición escrita y bajo su responsabilidad para que presenten, si lo creyesen conveniente el escrito alegando sobre el mérito de la prueba. Se considerará como una sola parte a quienes actúen bajo representación común. Transcurrido el plazo sin que el expediente haya sido devuelto, la parte que lo retuviese perderá el derecho de alegar sin que se requiera intimación. El plazo para presentar el alegato es común.
\end{formula}

\begin{note}
El texto vigente del art.\ 482 (según Ley 25.488) cierra con la frase «El plazo para presentar el alegato es común», mientras la docente describe el esquema sucesivo tradicional (seis días por parte, primero actores y luego demandados). Conviene verificar en cada juzgado cómo se computa el plazo en el sistema digital.
% TODO: contenido incompleto o ambiguo en las notas originales (plazo común o sucesivo)
\end{note}

\subsection{Visibilidad del alegato en el sistema digital}

El alegato se sube hoy digitalmente y permanece \textbf{secreto, no visible}, hasta que se cumpla el último plazo. Si el juzgado no lo pone visible a tiempo, la otra parte lo verá. Si se presenta muy adelantado, conviene pasar por el juzgado o presentar un escrito advirtiendo que se trata de alegatos y que no deben hacerse visibles. A veces se les pasa. La docente nunca vio que ocurriera, pero puede pasar.

\begin{important}
El \textbf{alegato} es el resumen final: \textit{qué probé, qué no pude probar, qué no probó la otra parte}. Para la sentencia (y para la Cámara) suele ser lo más leído. Plazo: \textbf{seis días por parte}; no se hace en familia ni en procesos especiales, salvo orden del juez.
\end{important}

\section{Cuadro Cornell: prueba y alegatos}

\begin{longtable}{L{0.27\textwidth} L{0.67\textwidth}}
\toprule
\textbf{Preguntas / palabras clave} & \textbf{Notas / respuestas} \\
\midrule
\endhead

\textbf{¿Cómo ofrezco la prueba?} &
Las pruebas son las taxativamente previstas en el Código; otras requieren autorización. Hay que explicar para qué se ofrece cada una: testigos (extremos a probar), informativa (qué debe contestar el informe), pericial (puntos de pericia). Los oficios a entidades privadas pueden ir directamente. \\
\addlinespace

\textbf{¿Cómo se abre a prueba? Audiencia 360 y 40 días} &
Con un proveído con fecha cierta; desde la audiencia corren \textbf{40 días hábiles} para producir la prueba. Suele haber conciliación previa; si no hay acuerdo, se pide la apertura a prueba y sale el proveído de prueba. \\
\addlinespace

\textbf{¿Qué es la caducidad o negligencia?} &
Se da cuando una parte no produce su prueba y deja vencer el plazo. Puede pedirse contra la otra parte, y la otra parte puede pedirla contra uno. Es una resolución apelable. \\
\addlinespace

\textbf{¿Qué hago si el proveído de prueba rechaza una prueba?} &
Es facultad ordenatoria del juez, inapelable en ese momento: se controla el proveído, se pide lo omitido y se reclama en la apelación de la sentencia (y también en los alegatos). Es distinto de la negligencia o caducidad, que sí se apelan. \\
\addlinespace

\textbf{¿Qué son los alegatos y cómo se estructuran?} &
Resumen final de la parte: qué probé, qué no pude probar y qué no probó la otra parte, y cómo me afectó la prueba que produjo la otra. Son lo más leído para sentenciar y para la Cámara. \\
\addlinespace

\textbf{¿Cuál es el plazo de los alegatos y cuándo no se hacen?} &
Seis días por parte (primero las actoras y luego las demandadas). No se hacen en familia ni en procesos especiales como el amparo, salvo orden del juez. El alegato digital queda oculto hasta vencer el último plazo. \\

\midrule
\multicolumn{2}{p{\dimexpr0.94\textwidth+2\tabcolsep\relax}}{\textbf{Resumen:} la prueba se ofrece de manera fundada, se abre con la audiencia 360 y se produce en 40 días hábiles. El proveído de prueba debe controlarse: el rechazo de una prueba es inapelable y se reclama recién en la apelación de la sentencia, a diferencia de la negligencia o caducidad, que sí se apelan. Cierran el proceso los alegatos (qué probé, qué no probé, qué no probó la otra parte), que son lo más leído por jueces y cámaras.} \\
\bottomrule
\end{longtable}

% ------------------------------------------------------------
% SÍNTESIS FINAL DE LA CLASE
% ------------------------------------------------------------
\chapter[Síntesis final]{Síntesis final de la clase}

La clase recorre el camino que sigue una demanda civil desde que se presenta hasta que las partes alegan, con un eje práctico: \textbf{cada acto vale solo si se notifica y se cuenta bien}.

La demanda y la sentencia van siempre por cédula (art.\ 135); los medios son taxativos (art.\ 136) y todo lo que los exceda, como WhatsApp o mail, exige autorización del juzgado, con mayor flexibilidad en familia. La cédula puede ser ley (extraña jurisdicción), papel o electrónica; se sigue con el libro de cédulas y el código de barras, y las demás resoluciones quedan notificadas por nota los martes y viernes, con plazos que las acordadas pueden alterar y que, ante la duda, conviene contar corto.

La palabra del proveído («notifíquese», «hágase saber», «traslado») dice cómo notificar; si no dice nada, suele ser un rechazo. Ampliar demanda o prueba exige rapidez y solo es posible antes de que la cédula se diligencie; después solo caben hechos nuevos.

La prueba se ofrece de manera fundada, se abre con la audiencia 360 y se produce en 40 días hábiles; el proveído de prueba debe controlarse porque su rechazo es inapelable y se reclama recién en la apelación de la sentencia, a diferencia de la negligencia o caducidad, que sí se apelan. Cierran el proceso los alegatos (qué probé, qué no probé, qué no probó la otra parte), que son lo más leído por jueces y cámaras.

\end{document}
